\PassOptionsToPackage{table,xcdraw}{xcolor}
\documentclass[10pt, a4paper, landscape]{article}

\usepackage[T1]{fontenc}
\usepackage[utf8]{inputenc}
\usepackage{tgpagella}
\usepackage[spanish, es-noshorthands]{babel}
\usepackage{amsmath, amssymb}
\usepackage{booktabs}
\usepackage{tabularx}
\usepackage{tikz}
\usepackage[american]{circuitikz}
\usetikzlibrary{arrows.meta, positioning, calc}
\usepackage[margin=1.5cm]{geometry}
\usepackage{enumitem}
\usepackage{microtype} % Optimiza la justificación y elimina micro-desbordamientos (Overfull \hbox)

\definecolor{ucbblue}{RGB}{0, 51, 102}
\definecolor{ucbgray}{RGB}{240, 244, 248}
\definecolor{warnred}{RGB}{180, 20, 20}

\begin{document}

\noindent
{\LARGE \textbf{\textcolor{ucbblue}{IMT-221 — Circuitos Electrónicos III}}}\\[0.15cm]
{\Large \textbf{Hito 3: Guía Única E1 $\rightarrow$ E4 — Parte 2: Instrumentación y Etapa E1 (NMOS Low-Side)}}\\[0.1cm]
{\small \textbf{Versión de trabajo 1.0} $\cdot$ 3 de noviembre de 2026 $\cdot$ Condicionado a firma de Puerta G1}
\par\vspace{0.2cm}\noindent\rule{\linewidth}{0.4pt}\vspace{0.3cm}

\section*{2. Esquema Común e Instrumentación}

\subsection*{2.1 Circuito E1 (Mantenido también como carga en E2)}

\begin{center}
\begin{circuitikz}[scale=0.95, transform shape, thick]
    % Alimentación y carga resistiva
    \draw (0,4.5) node[above]{\textbf{$+12\,\text{V}_{\text{BUS}}$} (Fuente limitada)} to[R, l={$R_L = 100\,\Omega$}, a={\small Potencia verificada}, i={$I_D$}] (0,2.0);

    % Transistor NMOS M1
    \draw (0,0.2) node[nigfete] (M1) {};
    \node[above left=0.05cm] at (M1.G) {G};
    \node[above right=0.05cm] at (M1.D) {D};
    \node[below right=0.05cm] at (M1.S) {S};
    \node[below right=0.05cm] at (M1.B) {B};

    \coordinate (drain) at ($(M1.D) + (0, 0.6)$);
    \coordinate (source) at ($(M1.S) + (0, -0.8)$);
    \coordinate (gate) at ($(M1.G) + (-2.2, 0)$);
    \coordinate (ctrl) at ($(gate) + (-3.0, 0)$);
    \coordinate (gnd_ctrl) at (ctrl |- source);
    \coordinate (gnd_bus) at ($(gnd_ctrl) + (0, -0.8)$);

    % Conexiones Drain, Source y Bulk
    \draw (0,2.0) -- (drain) node[circle, fill, inner sep=1.5pt]{} -- (M1.D);
    \draw (M1.S) -- (source) node[ground]{};
    \draw (M1.B) -- ++(0.6,0) |- (source);

    % Diodo intrínseco explícito a la derecha
    \draw[blue, dashed] (M1.S) -- ++(2.2,0) coordinate (db_s) to[D, l_={\small Diodo intrínseco}] (db_s |- M1.D) -- (M1.D);

    % Sondas de osciloscopio
    \draw[-{Stealth[length=2.5mm]}, warnred, very thick] (drain) -- ++(3.2,0) node[right]{\textbf{CH2}: Punta en D, ref.\ en S ($V_{DS}$)};
    \draw[-{Stealth[length=2.5mm]}, warnred, very thick] (gate) -- ++(0,1.6) node[above]{\textbf{CH1}: Punta en G, ref.\ en S ($V_{GS}$)};

    % Red de excitación de compuerta
    \draw (M1.G) -- (gate) node[circle, fill, inner sep=1.5pt]{};
    \draw (gate) to[R, l_={$R_G = 220\,\Omega$}] (ctrl) node[circle, fill, inner sep=1.5pt]{} node[above]{\textbf{Control}};
    \draw (ctrl) to[V, l_={\small $0-5\,\text{V}$ / $0-10\,\text{V}$}] (gnd_ctrl) node[ground]{};
    \draw (gate) to[R, l={$R_{PD} = 100\,\text{k}\Omega$}] (gate |- source) node[ground]{};

    % Retorno común verificado
    \draw[dashed, ucbblue] (gnd_bus) -- (source |- gnd_bus) node[midway, below]{\small Retorno común GND potencia / control (unir solo tras revisar tierras)};
\end{circuitikz}
\end{center}

\noindent
\textbf{Reglas de instrumentación física:}
\begin{itemize}[noitemsep, topsep=2pt]
    \item Revisar físicamente los terminales $G/D/S$ y la pestaña metálica del encapsulado TO-220 antes de cablear (la pestaña suele estar unida internamente al \textit{drain}).
    \item Identificar qué tierra está conectada a protección física; un osciloscopio de banco tiene los caimanes de referencia de sus sondas unidos a tierra de protección.
    \item \textbf{CH1:} punta en \textit{gate}, referencia en \textit{source} (mide $V_{GS}$ en low-side). \textbf{CH2:} punta en \textit{drain}, referencia en \textit{source} (mide $V_{DS}$). \textbf{\textcolor{warnred}{Prohibido conectar la pinza de tierra al drain ni al lado superior de la carga.}}
    \item Medir corriente mediante instrumento/método aprobado: un multímetro en modo amperímetro se conecta \textbf{en serie}, jamás en paralelo. Registrar escala, atenuación ($\times 1$ o $\times 10$), acoplamiento DC y límites.
\end{itemize}

\subsection*{2.2 Secuencia Universal Antes de una Energización (Todas las Etapas E1--E4)}
\begin{enumerate}[noitemsep, topsep=2pt]
    \item Dibujar el circuito y comprobar continuidad con \textbf{todas las fuentes apagadas}.
    \item Identificar componentes reales, pinout y límite de cada fuente, conductor y sonda.
    \item Preparar LTspice: guardar esquema, modelo/fuente, ventana de análisis transitorio y archivo \texttt{.log}.
    \item Verificar la \textbf{señal de control aislada de la potencia} en el osciloscopio ($V_{\min}$, $V_{\max}$, frecuencia $f$, ciclo de trabajo \textit{duty} y tiempos de flanco).
    \item Unir tierras \textbf{solo después} de revisar si existe conexión a tierra física y validar la topología.
    \item Comprobar que la compuerta tiene estado OFF definido al desconectar la señal ($R_{PD}$) y que jamás se excede $V_{GS(\max)}$.
    \item Presentar la puerta G correspondiente al docente; anotar \texttt{autorizado / observado / bloqueado}, fecha y firma.
    \item \textbf{Solo si está autorizado:} limitar corriente de fuente, comenzar en estado OFF, avanzar por estados aprobados, medir y detener ante calor, corriente o tensión fuera de rango.
    \item Apagar alimentación y descargar nodos antes de modificar cables. Guardar evidencia y discrepancias.
\end{enumerate}

\section*{3. Etapa E1 — Interruptor NMOS Low-Side Resistivo}

\subsection*{3.1 Derivación Matemática Paso a Paso}

\paragraph{1. Estado ON Ideal ($R_{DS(\text{on})} \rightarrow 0\,\Omega$, $V_{DS,\text{ON}}^{\text{ideal}} = 0\,\text{V}$):}
Aplicando la Ley de Tensiones de Kirchhoff (KVL) al lazo de potencia con $V_{\text{BUS}} = 12\,\text{V}$ y $R_L = 100\,\Omega$:
\begin{align}
    V_{\text{BUS}} - I_{D,\text{ON}}^{\text{ideal}} R_L - V_{DS,\text{ON}}^{\text{ideal}} &= 0 \implies I_{D,\text{ON}}^{\text{ideal}} = \frac{V_{\text{BUS}}}{R_L} = \frac{12\,\text{V}}{100\,\Omega} = 0.12\,\text{A} = 120\,\text{mA}. \\
    P_{RL,\text{ON}}^{\text{ideal}} &= \left(I_{D,\text{ON}}^{\text{ideal}}\right)^2 R_L = (0.12\,\text{A})^2 \times 100\,\Omega = 0.0144\,\text{A}^2 \times 100\,\Omega = 1.44\,\text{W}.
\end{align}

\paragraph{2. Estado OFF Ideal ($I_{D,\text{OFF}}^{\text{ideal}} = 0\,\text{A}$):}
\begin{equation}
    V_{DS,\text{OFF}}^{\text{ideal}} = V_{\text{BUS}} - I_{D,\text{OFF}}^{\text{ideal}} R_L = 12\,\text{V} - 0\,\text{V} = 12\,\text{V}, \qquad P_{RL,\text{OFF}}^{\text{ideal}} = 0\,\text{W}.
\end{equation}

\paragraph{3. Régimen Periódico PWM Ideal ($f = 100\,\text{Hz}$, $D = 0.50$):}
El periodo de conmutación es $T = \frac{1}{f} = \frac{1}{100\,\text{Hz}} = 0.010\,\text{s} = 10\,\text{ms}$, con $t_{\text{ON}} = D T = 5\,\text{ms}$ y $t_{\text{OFF}} = (1-D)T = 5\,\text{ms}$. La potencia media ideal disipada en la resistencia de carga es:
\begin{equation}
    P_{RL,\text{med}}^{\text{ideal}} = \frac{1}{T}\int_0^T p_{RL}(t)\,dt = \frac{1}{T}\int_0^{D T} P_{RL,\text{ON}}^{\text{ideal}}\,dt = D \cdot P_{RL,\text{ON}}^{\text{ideal}} = 0.50 \times 1.44\,\text{W} = 0.72\,\text{W}.
\end{equation}
\textit{Criterio térmico:} El resistor debe soportar tanto el promedio como estados ON continuos ($1.44\,\text{W}$); por tanto, un resistor de $1/4\,\text{W}$ ($0.25\,\text{W}$) es inadecuado (se sobrecarga en $1.44 / 0.25 = 5.76\times$).

\paragraph{4. Estado ON No Ideal (con $R_{DS(\text{on})}$ condicionado a $V_{GS}$, $I_D$ y $T_j$):}
\begin{align}
    V_{\text{BUS}} - I_D R_L - I_D R_{DS(\text{on})} &= 0 \implies I_D = \frac{V_{\text{BUS}}}{R_L + R_{DS(\text{on})}}, \\
    V_{DS,\text{on}} &= I_D R_{DS(\text{on})} = V_{\text{BUS}}\frac{R_{DS(\text{on})}}{R_L + R_{DS(\text{on})}}, \\
    P_{\text{MOS,cond}} &= I_D^2 R_{DS(\text{on})} = \left(\frac{V_{\text{BUS}}}{R_L + R_{DS(\text{on})}}\right)^2 R_{DS(\text{on})}.
\end{align}
Evaluando para un valor \textbf{hipotético y condicionado} $R_{DS(\text{on})} = 25\,\text{m}\Omega = 0.025\,\Omega$:
\begin{align*}
    I_D &= \frac{12\,\text{V}}{100\,\Omega + 0.025\,\Omega} = \frac{12\,\text{V}}{100.025\,\Omega} \approx 0.119970\,\text{A} = 119.97\,\text{mA}, \\
    V_{DS,\text{on}} &= 0.119970\,\text{A} \times 0.025\,\Omega \approx 0.002999\,\text{V} \approx 3.00\,\text{mV}, \\
    P_{\text{MOS,cond}} &= (0.119970\,\text{A})^2 \times 0.025\,\Omega \approx 0.0003598\,\text{W} \approx 0.36\,\text{mW}.
\end{align*}

\textbf{Advertencia instrumental:} No presentar $3\,\text{mV}$ como lectura garantizada, pues puede quedar por debajo del piso de ruido e incertidumbre del osciloscopio de aula. Asimismo, en el \textbf{IRLZ44N} el fabricante especifica conducción a niveles lógicos ($R_{DS(\text{on})}$ máximo a $4.5\,\text{V}$ o $5\,\text{V}$), mientras que en el \textbf{IRFZ44N} la especificación de $17.5\,\text{m}\Omega$ se garantiza únicamente a $V_{GS} = 10\,\text{V}$.

\subsection*{3.2 Señal de Control, LTspice E1 y Tabla de Referencia}
\begin{itemize}[noitemsep, topsep=2pt]
    \item \textbf{Generador principal ($100\,\text{Hz}$, cuadrada, $50\%$):} Para \texttt{IRLZ44N} buscar $0-5\,\text{V}$; para \texttt{IRFZ44N} buscar $0-10\,\text{V}$ ($10\,\text{V}_{\text{pp}}$ con $+5\,\text{V}$ de offset, verificado en osciloscopio sin gate y con gate). Si un generador \texttt{SKMI SK3000} entrega señal bipolar $-5/+5\,\text{V}$ sin offset ajustable, \textbf{prohibido conectarlo como si fuera $0-10\,\text{V}$}.
    \item \textbf{Arduino UNO $5\,\text{V}$ (opcional):} $100\,\text{Hz}$, duty $50\%$, con $R_G = 330\,\Omega$ inicial y $R_{PD} = 100\,\text{k}\Omega$:
\end{itemize}

\begin{verbatim}
const byte PIN_GATE = 9;
void setup() { digitalWrite(PIN_GATE, LOW); pinMode(PIN_GATE, OUTPUT); }
void loop()  { digitalWrite(PIN_GATE, HIGH); delay(5);
               digitalWrite(PIN_GATE, LOW);  delay(5); }
\end{verbatim}

\begin{itemize}[noitemsep, topsep=2pt]
    \item \textbf{LTspice E1:} Fuente DC \texttt{VBUS=12}, \texttt{RL=100}, control \texttt{PULSE(0 5 1m 1u 1u 5m 10m)} (o \texttt{0 10} para variante $10\,\text{V}$), \texttt{RG=220}, \texttt{RPD=100k}, directiva \texttt{.tran 0 31m 0 1u}. Si se usa \texttt{MDEMO}, rotularlo como \textbf{modelo ficticio}.
\end{itemize}

\par\vspace{0.15cm}\noindent
\renewcommand{\arraystretch}{1.25}
\begin{tabularx}{\linewidth}{|>{\raggedright\arraybackslash}X|c|c|c|c|}
\hline
\rowcolor{ucbgray}
\textbf{Estado (Referencia Ideal — NO Experimental)} & \textbf{$V_{GS}$ Ideal} & \textbf{$V_{DS}$ Ideal} & \textbf{$I_D$ Ideal} & \textbf{$P_{RL}$ Instantánea Ideal} \\ \hline
OFF & $0\,\text{V}$ & $12\,\text{V}$ & $0\,\text{A}$ & $0\,\text{W}$ \\ \hline
ON, \texttt{IRLZ44N} (Excitación lógica) & $5\,\text{V}$ & $\approx 0\,\text{V}$ & $\approx 120\,\text{mA}$ & $\approx 1.44\,\text{W}$ \\ \hline
ON, \texttt{IRFZ44N} (Excitación estándar) & $10\,\text{V}$ & $\approx 0\,\text{V}$ & $\approx 120\,\text{mA}$ & $\approx 1.44\,\text{W}$ \\ \hline
\end{tabularx}

\end{document}
